% Homework Assignment Article
% LaTeX Template
% Version 1.3.1 (ECL) (08/08/17)
%
% This template has been downloaded from:
% Overleaf
%
% Original author:
% Victor Zimmermann (zimmermann@cl.uni-heidelberg.de)
%
% License:
% CC BY-SA 4.0 (https://creativecommons.org/licenses/by-sa/4.0/)
%
%%%%%%%%%%%%%%%%%%%%%%%%%%%%%%%%%%%%%%%%%

%----------------------------------------------------------------------------------------

\documentclass[a4paper]{article} % Uses article class in A4 format

%----------------------------------------------------------------------------------------
%	FORMATTING
%----------------------------------------------------------------------------------------

\addtolength{\hoffset}{-2.25cm}
\addtolength{\textwidth}{4.5cm}
\addtolength{\voffset}{-3.25cm}
\addtolength{\textheight}{5cm}
\setlength{\parskip}{0pt}
\setlength{\parindent}{0in}

%----------------------------------------------------------------------------------------
%	PACKAGES AND OTHER DOCUMENT CONFIGURATIONS
%----------------------------------------------------------------------------------------

\usepackage[utf8]{inputenc} % Use UTF-8 encoding
\usepackage[T1]{fontenc}
\usepackage{lmodern}
\usepackage[english, spanish]{babel} % Language hyphenation and typographical rules

\usepackage{blindtext} % Package to generate dummy text
\usepackage{charter} % Use the Charter font
\usepackage{microtype} % Slightly tweak font spacing for aesthetics

\usepackage{amsthm, amsmath, amssymb} % Mathematical typesetting
\usepackage{float} % Improved interface for floating objects
\usepackage[final, colorlinks = true, 
            linkcolor = black, 
            citecolor = black]{hyperref} % For hyperlinks in the PDF
\usepackage{graphicx, multicol} % Enhanced support for graphics
\usepackage{xcolor} % Driver-independent color extensions
\usepackage{colortbl} % Para colores en tablas (Requerido por UML)
\usepackage{array} % Requerido por UML
\usepackage{marvosym, wasysym} % More symbols
\usepackage{rotating} % Rotation tools
\usepackage{censor} % Facilities for controlling restricted text
\usepackage{listings} % Environment for non-formatted code
\usepackage{pseudocode} % Environment for specifying algorithms in a natural way
\usepackage{booktabs} % Enhances quality of tables

% Configuración de listings
\lstset{
    language=Java,         
    basicstyle=\ttfamily\small,    
    keywordstyle=\color{blue!70!black}\bfseries,
    commentstyle=\color{green!50!black}\itshape,
    numbers=left,            
    frame=lines             
}

% Configuración TikZ y Árboles
\usepackage{tikz}
\usepackage{tikz-qtree} % Easy tree drawing tool
\usetikzlibrary{arrows.meta,positioning,fit,backgrounds,calc,shadows} % Librerías para UML
\tikzset{every tree node/.style={align=center,anchor=north},
         level distance=2cm} % Configuration for q-trees

\usepackage[backend=biber,style=numeric,
            sorting=nyt]{biblatex} % Complete reimplementation of bibliographic facilities
\usepackage{csquotes} % Context sensitive quotation facilities

\usepackage[yyyymmdd]{datetime} % Uses YEAR-MONTH-DAY format for dates
\renewcommand{\dateseparator}{-} % Sets dateseparator to '-'

\usepackage{fancyhdr} % Headers and footers
\pagestyle{fancy} % All pages have headers and footers
\fancyhead{}\renewcommand{\headrulewidth}{0pt} % Blank out the default header
\fancyfoot[L]{} % Custom footer text
\fancyfoot[C]{} % Custom footer text
\fancyfoot[R]{\thepage} % Custom footer text

\newcommand{\note}[1]{\marginpar{\scriptsize \textcolor{red}{#1}}} % Enables comments in red on margin

%----------------------------------------------------------------------------------------
%	UML DIAGRAM CONFIGURATION
%----------------------------------------------------------------------------------------

% --- Color palette ---
\definecolor{umlHeader}{RGB}{215,232,252}
\definecolor{umlHeader2}{RGB}{226,239,218}
\definecolor{umlBorder}{RGB}{54,84,134}
\definecolor{namespaceFill}{RGB}{248,250,253}
\definecolor{assocColor}{RGB}{70,70,70}
\definecolor{inheritColor}{RGB}{45,78,120}

\newcolumntype{L}[1]{>{\raggedright\arraybackslash}p{#1}}

% UML class macro: #1 node name, #2 position, #3 width, #4 header color, #5 class name, #6 attributes, #7 methods
\newcommand{\umlclass}[7]{%
  \node[inner sep=0pt, outer sep=0pt, drop shadow={opacity=0.25, shadow xshift=2pt, shadow yshift=-2pt}] (#1) at #2 {%
    \scriptsize\sffamily
    \renewcommand{\arraystretch}{1.10}%
    \setlength{\tabcolsep}{4pt}%
    \begin{tabular}{|L{#3}|}
      \hline
      \rowcolor{#4}\multicolumn{1}{|c|}{\bfseries #5}\\
      \hline
      #6\\
      \hline
      #7\\
      \hline
    \end{tabular}%
  };%
}

\begin{document}

%----------------------------------------------------------------------------------------
%	TITLE SECTION
%----------------------------------------------------------------------------------------

\fancyhead[C]{}
\hrule \medskip % Upper rule
\begin{minipage}{0.295\textwidth} % Left side of title section
\raggedright
UNAL\\ % Your lecture or course
\footnotesize % Authors text size
\hfill\\
Cristian Ruiz Hernandez\\
Mateo Quiceno Zapata
\end{minipage}
\begin{minipage}{0.4\textwidth} % Center of title section
\centering 
\large % Title text size
Actividad 5\\ % Assignment title and number
\normalsize % Subtitle text size
Programación Orientada a Objetos (3007744)\\ % Assignment subtitle
\end{minipage}
\begin{minipage}{0.295\textwidth} % Right side of title section
\raggedleft
\today\\ % Date
\footnotesize % Email text size
\hfill\\
cruizh@unal.edu.co\\
mquicenoz@unal.edu.co
\end{minipage}
\medskip\hrule % Lower rule
\bigskip

%----------------------------------------------------------------------------------------
%	ARTICLE CONTENTS
%----------------------------------------------------------------------------------------

\section{Biblioteca CRUD (Persistencia en Archivos)}
Desarrollo de una Aplicación con Interfaz Gráfica (GUI) conectada a un archivo de datos local. El sistema debe contener un formulario provisto de botones funcionales para ejecutar de manera exitosa el patrón de persistencia \textbf{CRUD} (Create, Read, Update y Delete) sobre un archivo de texto secuencial usando Java. Las operaciones de persistencia deben interactuar con el sistema de archivos del sistema operativo a través del uso de la clase \texttt{RandomAccessFile}, gestionando la integridad, evitando duplicados, y permitiendo la actualización y eliminación física de los registros.

\bigskip

\subsection{Interfaz de usuario}
A continuación se detallan los componentes y estados de la interfaz gráfica desarrollada para la gestión y administración de la biblioteca personal:

\begin{enumerate}
    \item \textbf{Ventana de Inicio (Limpia):} Estado inicial de la aplicación mostrando un formulario vacío y una lista sin elementos en la tabla, indicando una base de datos limpia.
    \begin{figure}[H]
        \centering
        \includegraphics[width=0.45\textwidth]{imagenes/1.png}
        \caption{Interfaz inicial limpia para la administración de libros.}
    \end{figure}

    \item \textbf{Diligenciamiento de Datos:} El usuario completa la información (Título, Autor, Estado) en el formulario de la interfaz antes de presionar el botón de creación.
    \begin{figure}[H]
        \centering
        \includegraphics[width=0.45\textwidth]{imagenes/2.png}
        \caption{Datos cargados en el formulario de registro.}
    \end{figure}

    \item \textbf{Confirmación de Creación Exitosa:} Tras hacer clic en ``Añadir (Create)'', se escribe en el archivo físico y se despliega una alerta de confirmación.
    \begin{figure}[H]
        \centering
        \includegraphics[width=0.45\textwidth]{imagenes/3.png}
        \caption{Diálogo emergente informando el guardado exitoso del nuevo registro.}
    \end{figure}

    \item \textbf{Tabla con Registro Visualizado (Read):} Al cerrarse el diálogo, la tabla actualiza su listado realizando una lectura completa del archivo secuencial.
    \begin{figure}[H]
        \centering
        \includegraphics[width=0.45\textwidth]{imagenes/4.png}
        \caption{Tabla que refleja el registro recién creado en el archivo \texttt{biblioteca.txt}.}
    \end{figure}

    \item \textbf{Selección y Carga en Formulario:} Al seleccionar una fila de la tabla, los campos del formulario se rellenan automáticamente permitiendo al usuario realizar modificaciones sobre el objeto.
    \begin{figure}[H]
        \centering
        \includegraphics[width=0.45\textwidth]{imagenes/5.png}
        \caption{Fila seleccionada, transfiriendo la información correspondiente a los inputs de texto.}
    \end{figure}

    \item \textbf{Mensaje de Edición Exitosa (Update):} Tras presionar ``Guardar Edición (Update)'', el sistema sobreescribe el archivo mediante un archivo temporal y reestablece los campos.
    \begin{figure}[H]
        \centering
        \includegraphics[width=0.45\textwidth]{imagenes/6.png}
        \caption{Diálogo informativo de confirmación del proceso de edición física del registro.}
    \end{figure}

    \item \textbf{Control de Campos Nulos o Vacíos:} El sistema valida que los campos de Título y Autor contengan texto válido. Si están vacíos, detiene la operación y alerta al usuario.
    \begin{figure}[H]
        \centering
        \includegraphics[width=0.45\textwidth]{imagenes/7.png}
        \caption{Mensaje restrictivo impidiendo el registro de datos incompletos.}
    \end{figure}

    \item \textbf{Evitar Duplicidad de Registros:} Al intentar crear un libro cuyo título ya existe en el archivo, se levanta una alerta impidiendo el guardado repetido.
    \begin{figure}[H]
        \centering
        \includegraphics[width=0.45\textwidth]{imagenes/8.png}
        \caption{Excepción controlada por pantalla previniendo redundancia en la persistencia.}
    \end{figure}

    \item \textbf{Confirmación de Eliminación (Delete):} Al presionar ``Eliminar (Delete)'', se solicita confirmación de exclusión al usuario antes de borrar físicamente el registro en el archivo.
    \begin{figure}[H]
        \centering
        \includegraphics[width=0.45\textwidth]{imagenes/9.png}
        \caption{Diálogo de confirmación del tipo Yes/No para proteger la consistencia de los datos.}
    \end{figure}
\end{enumerate}

\subsection{Diagrama de clases}

\begin{center}
\resizebox{0.85\textwidth}{!}{%
\begin{tikzpicture}[
  font=\sffamily,
  inheritance/.style={draw=inheritColor, line width=0.8pt, -{Triangle[open, length=6pt, width=8pt]}},
  association/.style={draw=assocColor, line width=0.65pt, -{Latex[length=4pt, width=5pt]}},
  plainassoc/.style={draw=assocColor, line width=0.65pt},
  label/.style={font=\scriptsize\sffamily, fill=white, inner sep=1.5pt},
  multiplicity/.style={font=\scriptsize\sffamily, fill=white, inner sep=1pt}
]

% -----------------------------------------------------------------------------
% Classes in namespace Actividad5
% -----------------------------------------------------------------------------

\umlclass{VentanaBiblioteca}{(0,0)}{5.5cm}{umlHeader}{VentanaBiblioteca}{%
-txtTitulo: JTextField\\
-txtAutor: JTextField\\
-cbEstado: JComboBox\\
-tablaLibros: JTable\\
-modeloTabla: DefaultTableModel\\
-tituloSeleccionado: String}{%
+VentanaBiblioteca()\\
-limpiarFormulario(): void\\
-cargarTabla(): void\\
-crear(): void\\
-actualizar(): void\\
-borrar(): void\\
+main(String[] args)\$}

\umlclass{GestorArchivoBiblioteca}{(7.5,0)}{5.2cm}{umlHeader}{GestorArchivoBiblioteca}{%
-FILE\_NAME: String = ``biblioteca.txt''}{%
+leerLibros(): ArrayList\$\\
+crearLibro(String, String, String)\$\\
+actualizarLibro(String, String, String, String)\$\\
+borrarLibro(String)\$}

\umlclass{UIUtils}{(0,-6.5)}{5.0cm}{umlHeader}{UIUtils}{%
+COLOR\_FONDO: Color\$\\
+COLOR\_PRIMARIO: Color\$\\
+COLOR\_TEXTO: Color\$\\
+FUENTE\_TITULO: Font\$\\
+FUENTE\_NORMAL: Font\$}{%
+aplicarTema()\$\\
+estilizarBoton(JButton)\$\\
+estilizarPanel(JPanel)\$}

% -----------------------------------------------------------------------------
% Namespace boundary Actividad5
% -----------------------------------------------------------------------------
\begin{scope}[on background layer]
  \node[
    draw=umlBorder,
    rounded corners=8pt,
    very thick,
    fill=namespaceFill,
    inner xsep=14pt,
    inner ysep=26pt,
    fit=(VentanaBiblioteca) (GestorArchivoBiblioteca) (UIUtils)
  ] (namespace) {};
  \node[
    anchor=north west,
    fill=umlBorder,
    text=white,
    rounded corners=3pt,
    font=\bfseries\sffamily\small,
    inner xsep=8pt,
    inner ysep=4pt
  ] at ([xshift=10pt,yshift=-8pt]namespace.north west) {namespace Actividad5};
\end{scope}

% -----------------------------------------------------------------------------
% Associations Actividad5
% -----------------------------------------------------------------------------
\draw[association] (VentanaBiblioteca.east) -- (GestorArchivoBiblioteca.west)
  node[label, midway, above] {Usa}
  node[multiplicity, near start, above] {1}
  node[multiplicity, near end, above] {1};

\draw[association] (VentanaBiblioteca.south) -- (UIUtils.north)
  node[label, midway, right] {Aplica Estilo}
  node[multiplicity, near start, left] {1}
  node[multiplicity, near end, left] {1};

\end{tikzpicture}%
}
\end{center}

\subsection{Casos de uso}
\begin{itemize}
    \item \textbf{Caso de Uso 1: Crear Registro de Libro}
    \begin{itemize}
        \item \textbf{Actor:} Usuario / Bibliotecario.
        \item \textbf{Precondición:} La ventana principal de la aplicación debe estar desplegada.
        \item \textbf{Flujo Principal:} El usuario escribe el Título, el Autor y selecciona el Estado. Presiona el botón \textit{Añadir (Create)}. El sistema valida que los campos no estén vacíos y que el título no sea redundante, escribe el registro en el archivo de texto y refresca la tabla visual.
        \item \textbf{Flujo Alternativo (Campos Vacíos):} Si algún campo obligatorio está vacío, el sistema interrumpe la lógica y notifica mediante una alerta emergente de tipo advertencia.
        \item \textbf{Flujo Alternativo (Duplicidad):} Si el título ya existe en el archivo, se arroja una alerta de error y se bloquea la escritura.
    \end{itemize}
    
    \item \textbf{Caso de Uso 2: Leer Libros Almacenados}
    \begin{itemize}
        \item \textbf{Actor:} Sistema / Usuario.
        \item \textbf{Precondición:} El archivo \texttt{biblioteca.txt} debe existir o crearse en el primer acceso.
        \item \textbf{Flujo Principal:} Al inicializar la interfaz o tras cualquier cambio, el sistema ejecuta la lectura del archivo plano línea por línea. Decodifica el formato delimitado por el caracter especial \texttt{!} en variables y las renderiza ordenadamente dentro del modelo de datos de la \texttt{JTable}.
    \end{itemize}

    \item \textbf{Caso de Uso 3: Actualizar Estado del Libro}
    \begin{itemize}
        \item \textbf{Actor:} Usuario.
        \item \textbf{Precondición:} Un registro debe seleccionarse de la tabla cargada.
        \item \textbf{Flujo Principal:} El usuario selecciona un libro de la lista, modifica su estado o información en el formulario y pulsa \textit{Guardar Edición (Update)}. El sistema lee secuencialmente el archivo copiando las líneas en un archivo temporal; cuando detecta el título coincidente, inyecta la nueva línea modificada. Reemplaza el archivo temporal por el original y actualiza la vista.
    \end{itemize}

    \item \textbf{Caso de Uso 4: Eliminar Libro}
    \begin{itemize}
        \item \textbf{Actor:} Usuario.
        \item \textbf{Precondición:} Un registro debe estar previamente seleccionado.
        \item \textbf{Flujo Principal:} El usuario pulsa \textit{Eliminar (Delete)}. El sistema despliega un mensaje confirmando la acción (Yes/No). Si el usuario selecciona sí, el sistema realiza una copia de exclusión secuencial (omitiendo la línea del libro seleccionado) mediante un archivo temporal, sobreescribe el archivo principal, limpia el formulario y actualiza la lista.
    \end{itemize}
\end{itemize}

\subsection{Link}
Para verificar funcionamiento se debe ejecutar el archivo de "VentanaBiblioteca.java"\\
\url{https://github.com/mateolikescats/POO/tree/main/Actividad_5/src/Actividad5}

\end{document}
